\textbf{Phần 1. Trắc nghiệm lựa chọn}
\setcounterpageref{ex}{0}

\begin{ex} %[2D1N04]
Muốn biết số thứ nhất hơn số thứ hai bao nhiêu đơn vị, ta làm phép tính gì?
\choice
{Phép cộng}
{Phép nhân}
{\True Phép trừ (lấy số lớn trừ số bé)}
{Phép chia}
\loigiai{
Ta lấy số lớn trừ số bé.\\
\textbf{Đáp án của câu hỏi là C}
}
\end{ex}

\begin{ex} %[2D1N04]
Số 15 hơn số 10 bao nhiêu đơn vị?
\choice
{$3$}
{\True $5$}
{$25$}
{$4$}
\loigiai{
$15 - 10 = 5$.\\
\textbf{Đáp án của câu hỏi là B}
}
\end{ex}

\begin{ex} %[2D1H04]
Số 20 kém số 28 bao nhiêu đơn vị?
\choice
{\True $8$}
{$6$}
{$10$}
{$48$}
\loigiai{
$28 - 20 = 8$.\\
\textbf{Đáp án của câu hỏi là A}
}
\end{ex}

\begin{ex} %[2D1H04]
Bạn Nam có 18 viên bi, bạn Hải có 12 viên bi. Nam có nhiều hơn Hải số viên bi là
\choice
{$5$ viên}
{\True $6$ viên}
{$7$ viên}
{$30$ viên}
\loigiai{
$18 - 12 = 6$ viên bi.\\
\textbf{Đáp án của câu hỏi là B}
}
\end{ex}

\begin{ex} %[2D1H04]
Cây cam có 35 quả, cây bưởi có 20 quả. Cây bưởi có ít hơn cây cam số quả là
\choice
{\True $15$ quả}
{$10$ quả}
{$25$ quả}
{$55$ quả}
\loigiai{
$35 - 20 = 15$ quả.\\
\textbf{Đáp án của câu hỏi là A}
}
\end{ex}

\begin{ex} %[2D1H04]
Một đoạn dây dài $48\text{ cm}$, một đoạn dây khác dài $30\text{ cm}$. Đoạn thứ nhất dài hơn đoạn thứ hai là
\choice
{$16\text{ cm}$}
{$20\text{ cm}$}
{\True $18\text{ cm}$}
{$78\text{ cm}$}
\loigiai{
$48 - 30 = 18\text{ cm}$.\\
\textbf{Đáp án của câu hỏi là C}
}
\end{ex}

\begin{ex} %[2D1V04]
Tổ 1 trồng được 25 cây, tổ 2 trồng được 15 cây. Tổ 1 trồng nhiều hơn tổ 2 số cây là
\choice
{\True $10$ cây}
{$8$ cây}
{$12$ cây}
{$40$ cây}
\loigiai{
$25 - 15 = 10$ cây.\\
\textbf{Đáp án của câu hỏi là A}
}
\end{ex}

\begin{ex} %[2D1V04]
Bình có 24 cái kẹo, Lan có 36 cái kẹo. Bình có ít hơn Lan bao nhiêu cái kẹo?
\choice
{$10$ cái}
{$14$ cái}
{$60$ cái}
{\True $12$ cái}
\loigiai{
$36 - 24 = 12$ cái kẹo.\\
\textbf{Đáp án của câu hỏi là D}
}
\end{ex}

\begin{ex} %[2D1V04]
Số lớn nhất có hai chữ số hơn số bé nhất có hai chữ số bao nhiêu đơn vị?
\choice
{$80$}
{\True $89$}
{$90$}
{$88$}
\loigiai{
$99 - 10 = 89$.\\
\textbf{Đáp án của câu hỏi là B}
}
\end{ex}

\begin{ex} %[2D1V04]
Một lớp học có 20 bạn nam và 18 bạn nữ. Số bạn nam nhiều hơn số bạn nữ là
\choice
{$1$ bạn}
{$3$ bạn}
{$4$ bạn}
{\True $2$ bạn}
\loigiai{
$20 - 18 = 2$ bạn.\\
\textbf{Đáp án của câu hỏi là D}
}
\end{ex}

\begin{ex} %[2D1V04]
Bao gạo to nặng $50\text{ kg}$, bao gạo nhỏ nặng $20\text{ kg}$. Bao to nặng hơn bao nhỏ bao nhiêu ki-lô-gam?
\choice
{$25\text{ kg}$}
{$35\text{ kg}$}
{\True $30\text{ kg}$}
{$70\text{ kg}$}
\loigiai{
$50 - 20 = 30\text{ kg}$.\\
\textbf{Đáp án của câu hỏi là C}
}
\end{ex}

\begin{ex} %[2D1C04]
An có 15 viên bi, Bình có 25 viên bi. Hỏi Bình phải cho An bao nhiêu viên bi để số bi của hai bạn bằng nhau?
\choice
{$10$ viên}
{$6$ viên}
{$4$ viên}
{\True $5$ viên}
\loigiai{
Bình hơn An: $25 - 15 = 10$ viên bi.\\
Bình cần cho An: $10 : 2 = 5$ viên bi.\\
\textbf{Đáp án của câu hỏi là D}
}
\end{ex}

\textbf{Phần 2. Trắc nghiệm đúng sai}
\setcounterpageref{ex}{0}

\begin{ex} %[2D1H04]
Xét tính đúng sai về so sánh hơn kém:
\choiceTF
{\True Số 45 hơn số 30 là 15 đơn vị.}
{\True Số 20 kém số 50 là 30 đơn vị.}
{Số 60 hơn số 40 là 30 đơn vị.}
{\True Số 60 hơn số 40 là 20 đơn vị ($60 - 40 = 20$).}
\loigiai{
\begin{itemize}
    \item $45 - 30 = 15$: Đúng.
    \item $50 - 20 = 30$: Đúng.
    \item $60 - 40 = 20 \neq 30$. Ý 3 Sai, Ý 4 Đúng.
\end{itemize}
\begin{itemize}
    \item a) \textbf{Đúng}.
    \item b) \textbf{Đúng}.
    \item c) \textbf{Sai}.
    \item d) \textbf{Đúng}.
\end{itemize}
}
\end{ex}

\begin{ex} %[2D1V04]
Mai có 28 bông hoa, Cúc có 16 bông hoa. Xét tính đúng sai:
\choiceTF
{\True Mai có nhiều hơn Cúc 12 bông hoa ($28 - 16 = 12$).}
{\True Cúc có ít hơn Mai 12 bông hoa.}
{Cả hai bạn có tất cả 40 bông hoa.}
{\True Cả hai bạn có tất cả 44 bông hoa ($28 + 16 = 44$).}
\loigiai{
\begin{itemize}
    \item Mai hơn Cúc: $28 - 16 = 12$: Đúng.
    \item Cúc ít hơn Mai 12: Đúng.
    \item Tổng: $28 + 16 = 44 \neq 40$. Ý 3 Sai, Ý 4 Đúng.
\end{itemize}
\begin{itemize}
    \item a) \textbf{Đúng}.
    \item b) \textbf{Đúng}.
    \item c) \textbf{Sai}.
    \item d) \textbf{Đúng}.
\end{itemize}
}
\end{ex}

\textbf{Phần 3. Trả lời ngắn}
\setcounterpageref{ex}{0}

\begin{ex} %[2D1H04]
Số 58 hơn số 25 bao nhiêu đơn vị?
\shortans{33}
\loigiai{
\textbf{Đáp số:} $33$\\
$58 - 25 = 33$.
}
\end{ex}

\begin{ex} %[2D1H04]
Số 32 kém số 67 bao nhiêu đơn vị?
\shortans{35}
\loigiai{
\textbf{Đáp số:} $35$\\
$67 - 32 = 35$.
}
\end{ex}

\begin{ex} %[2D1V04]
Bác Ba nuôi 48 con gà và 25 con vịt. Hỏi số gà nhiều hơn số vịt bao nhiêu con?
\shortans{23}
\loigiai{
\textbf{Đáp số:} $23$\\
$48 - 25 = 23$ con.
}
\end{ex}

\begin{ex} %[2D1V04]
Hùng có 35 viên bi, Dũng có 14 viên bi. Hỏi Dũng có ít hơn Hùng bao nhiêu viên bi?
\shortans{21}
\loigiai{
\textbf{Đáp số:} $21$\\
$35 - 14 = 21$ viên.
}
\end{ex}

\begin{ex} %[2D1V04]
Số tròn chục lớn nhất có hai chữ số hơn số tròn chục bé nhất có hai chữ số bao nhiêu đơn vị?
\shortans{80}
\loigiai{
\textbf{Đáp số:} $80$\\
$90 - 10 = 80$.
}
\end{ex}

\begin{ex} %[2D1C04]
Mẹ 36 tuổi, con 12 tuổi. Mẹ hơn con bao nhiêu tuổi?
\shortans{24}
\loigiai{
\textbf{Đáp số:} $24$\\
$36 - 12 = 24$ tuổi.
}
\end{ex}
